Quisiera agradecer a la Universidad Nacional de Córdoba y a la Facultad de Matemática, Astronomía, Física y Computación por haberme dado la posibilidad de acceder a una educación libre, gratuita y de calidad, y por permitirme realizar este trabajo final de grado que culmina una etapa importante en mi vida y da comienzo a otra.

A mis directores, Damián Barsotti, Gabriel Della Bella y Pablo Barttfeld, por su guía, su paciencia y su compromiso en todo momento.

A cada compañero y profesor con quienes tuve la oportunidad de compartir estos años de estudio, por la compañía y por todo el conocimiento y las enseñanzas que me brindaron.

A mis compañeros de trabajo en Ualabee por permitirme desarrollarme profesionalmente, por las oportunidades que me brindaron y por darme el espacio para enfocarme en mis estudios cuando lo necesitaba.

A mis amigos de la facultad, por haberme enseñado tanto en lo académico como en lo personal, por las anécdotas, las horas de estudio y por hacer de la vida universitaria algo más divertido.
A mis amigos fuera de la facultad, que también me acompañaron y estuvieron dispuestos a escucharme y ayudarme cuando lo necesitaba.

Por último y más importante, a mi familia quienes son mi sostén y apoyo incondicional.

A mi papá Marcelo y a mi mamá Judith, quienes se esforzaron por darme la mejor educación posible y me brindaron todo para que no me faltara nada, incluso en los momentos más difíciles.
A mis hermanos, Matías y Carolina, que se han preocupado por mí, aconsejado y acompañado en todo momento.
A mis sobrinos, Felipe y Alma, que siempre me sacan una sonrisa y traen alegría a la familia.
A mis tíos y tías, abuelos y abuelas; en particular a mis tías Marta y Silvia, que se han preocupado por mi educación y mi salud en todo momento.
A todos ellos, a los que están y a los que ya no están con nosotros pero siguen presentes en el recuerdo, por creer en mí incluso cuando yo no lo hacía y por darme la fuerza para seguir adelante.

Y por último, a mí mismo. Por todo el esfuerzo, el sacrificio, la paciencia y la dedicación que puse en estos años de carrera y formación, a pesar de las dificultades que se presentaron en el camino. Por haber aprendido a enfrentarlas y superarlas.
